% !PRINTED 214
\weberpage{214}
%<<CONT
Die Gleichung (9), die in Bezug auf $\tau$ identisch ist, bleibt
aber auch richtig, wenn nach Ausf\"uhrung der angedeuteten
Potenzirung rechts $\tau^k$ durch $\tau_k$ ersetzt wird, und sie lehrt also,
dass zwischen den Functionen $Y(t, x)$, $H(\tau, \xi)$ die Relation besteht

\weq{12}{H(\tau, \xi) = r\, Y(t, x),}

wenn wir die $t$ von den $\tau$ in folgender Weise abh\"angen lassen:
Man setze

\weq{13}{\begin{aligned}
t_{n-2} &= (\alpha\tau+\beta)^{n-2}\\
t_{n-3} &= (\alpha\tau+\beta)^{n-3}(\gamma\tau+\delta)\\
&\ \cdots\cdots\cdots\cdots\cdots\cdots\cdots\cdots\cdots\\
t_0 &= (\gamma\tau+\delta)^{n-2},
\end{aligned}}

und ersetze \gesperrt{nach Ausf\"uhrung der Potenzen} $\tau^k$ durch $\tau_k$.

Es sei nun $z$ eine beliebige Ver\"anderliche, und wir multipliciren,
ehe wir die Potenzirungen in (13) ausf\"uhren, diese
Gleichungen der Reihe nach mit

\[
1, \quad -(n-2)z, \quad \frac{(n-2)(n-3)}{1\,.\,2}z^2, \cdots \pm z^{n-2};
\]

dann bekommen wir links eine ganze rationale Function von $z$
vom $(n-2)^{\text{ten}}$ Grade

\weq{14}{\begin{aligned}
T(z) &= t_{n-2} - (n-2)t_{n-3}z\\
&\quad + \frac{(n-2)(n-3)}{1\,.\,2}t_{n-4}z^2 - \cdots \pm t_0 z^{n-2},
\end{aligned}}

und die rechte Seite ergiebt nach dem binomischen Satz

\[
[(\alpha\tau+\beta) - z(\gamma\tau+\delta)]^{n-2} = [\tau(\alpha-\gamma z) + (\beta-\delta z)]^{n-2},
\]

was, wenn wir

\weq{15}{z = \frac{\alpha\xi+\beta}{\gamma\xi+\delta}, \qquad \xi = \frac{\delta z-\beta}{-\gamma z+\alpha}}

setzen, in

\[
(\alpha-\gamma z)^{n-2}(\tau-\xi)^{n-2} = \frac{r^{n-2}(\tau-\xi)^{n-2}}{(\gamma\xi+\delta)^{n-2}}
\]

\"ubergeht.

Wenn wir hierin die Potenz $(\tau-\xi)^{n-2}$ nach dem binomischen
Satze ausf\"uhren, und dann $\tau^k$ durch $\tau_k$ ersetzen, so erhalten
wir eine Umformung der Function $T(z)$. Wir setzen also

% [sic] Im Druck steht vor \cdots ein Pluszeichen; die Binomialentwicklung und (14) verlangen dort ein Minuszeichen.

\[
\begin{aligned}
\Theta(\xi) &= \tau_{n-2} - (n-2)\tau_{n-3}\xi\\
&\quad + \frac{(n-2)(n-3)}{1\,.\,2}\tau_{n-4}\xi^2 + \cdots \pm \tau_0\xi^{n-2}
\end{aligned}
\]

und erhalten die identische Umformung
%CONT>>
